% Section 9: DiD material
\section{Changes across periods and difference-in-differences}\label{sec:did}

Only one DiD result has been through external review: PHASE2b Proposition~D8,
which treats election periods as ``districts'' under M1. The rest of the DiD
material, in PHASE1b (\texttt{PHASE1b\_monotone\_did.md}) and PHASE2d
(\texttt{PHASE2d\_did.md}), is a \textbf{proposal that has not been reviewed}
(\texttt{MSIGN\_REVIEW\_LOG.md} item M-13; \texttt{EXTERNAL\_REVIEW\_PLAN.md}).
It is summarised here only so that the status table is complete; it is not
reproduced in the appendices, and nothing in Sections~\ref{sec:rxc}--\ref{sec:inference}
depends on it.

\subsection{The reviewed result: Proposition D8}

\begin{result}[Proposition D8, the change in rates under M1 with periods as districts --- PHASE2b \S5]\label{res:d8}
\status{Proved.} Let $D\in\{\text{pre},\text{post}\}$; assume M1, (S1)--(S2-d),
$P(X\in A_D)=1$ and A5-M1. With per-period moments $\mu^{[t]}$,
\[
B^*_{rc,\text{post}}-B^*_{rc,\text{pre}}=\text{(identified)}+L(K),\qquad
L(K)=\sum_{s\ne r}K_{c,rs}\,\Delta q_{rs},\quad
\Delta q_{rs}=\frac{\mu^{[\text{post}]}_{rs}}{\mu^{[\text{post}]}_r}-\frac{\mu^{[\text{pre}]}_{rs}}{\mu^{[\text{pre}]}_r}.
\]
(a) \emph{Identification by cancellation:} the change does not depend on $K$ for
every $K\in\Kspace$ iff $\Delta q_{rs}=0$ for all $s\ne r$ (for example when the
size-and-composition moments are equal across periods). (b) \emph{Identification
given the feasible set:} the change is point-identified iff
$\max_{K\in\Theta}L(K)-\min_{K\in\Theta}L(K)=0$, where $\Theta$ is the identified
set for $K$. This holds under (a), and also when $\Theta$ fixes $L$ although
$\Delta q\ne0$ (example: $2\times2$, $Q_{1,12}=0$, LCE on the whole simplex, so
$\Theta=\{0\}$).
\end{result}
\begin{plain}
If contextual effects are unchanged between two elections and only baselines move
(M1 across periods), the change in ecological transition rates is pinned down when
the electorate's size-weighted composition moments are the same in both periods,
or when the bounds happen to fix the unknown combination. Otherwise the change is
bounded by projecting the common identified set. This is a statement about the
difference of two ecological rates. A causal (difference-in-differences) reading
needs further structure, developed in PHASE2d, which is not yet reviewed.
\end{plain}
\proofref{Appendix~\ref{app:phase2b}, \S5, Proposition~D8 steps 1--5; check D8b [F]}
\begin{revnote}
(R-24; A2-4.) Round~1: promoted from Remark~6.2 to a proposition with the exact
moment condition. Round~2: the proposition said the change is point-identified,
whatever the identified set for $K$, \emph{iff} $\Delta q=0$; that is false as a
statement about identification (referee's example). It now separates (a)
cancellation, which is sufficient, from (b), necessary and sufficient given the
identified set, and drops the causal reading.
\end{revnote}

\subsection{Proposals not yet reviewed (PHASE1b, PHASE2d)}

PHASE2d develops DiD bounds for the Wisconsin and Iowa suffrage referenda (rows:
veterans and non-veterans; columns: yes, no, abstain). Its results are written
with proofs and checked in \texttt{verify\_phase2d.py} (16 PASS) and
\texttt{verify\_phase1\_mei\_did.py}, but they have \textbf{not been reviewed}.
In brief:
\begin{itemize}[nosep]
\item \textbf{Lemma G}: a general sharpness lemma for affine mean models with
convex selection sets (as PHASE2b Thm~1$'$).
\item \textbf{Proposition $\Delta0$}: with no link between periods, the DiD of the
gap lies in the difference of the two periods' DD ranges; widths add.
\item \textbf{Lemma FD1}: under first-difference LCE the regression of $\Delta T$
identifies $\Delta Q$ and leaves $\Delta K$ unidentified.
\item \textbf{Propositions FD2, FD3 and FD4}: single restrictions R1 (no differential
spillover onto non-veterans), R2 and R4 (Goodman on changes) each point-identify
the DiD, with R4 between R1 and R2; a bounded change $|\Delta w_1|\le\delta$ gives
an interval and a breakdown value; FD4 is the outer set with no restriction on
$\Delta K$.
\item \textbf{Proposition $\Delta4$}: Monotone EI applied to individual changes;
the DiD lies between Goodman-on-changes and the neighbourhood model on changes
under CR on changes.
\item \textbf{Proposition $\Delta6$}: with changing population counts, R1 gives a
short interval rather than a point.
\item \textbf{Propositions C1--C2, Remark C3}: composition change (stayers,
aging-in, exits).
\item \textbf{Proposition CL}: monotone treatment response, no decline absent
service and a selection assumption bound the ATT and are jointly refutable.
\item \textbf{Propositions P2--P3}: party information (a structural zero alone does
not help; selection on party only does).
\end{itemize}
PACKAGE\_DESIGN \S7 lists the change-versions of the modules (M-SIGN-$\Delta$,
M-ORD-$\Delta$, M-LOC-$\Delta$ and the sandwich DiD-S); DiD-S validity is
proved there by applying P-S to changes, and the rest of R11 is open. The DiD
rungs were renamed $\Delta0$--$\Delta8$ to avoid a clash with PHASE2b's
Propositions~D0--D8 (M-14). The PHASE1b and PHASE2d clustering language was
revised to ``one unit observed twice'' and to county effects in the mean (M-5).
